Particle correlations underpin a wide range of calculations in molecular dynamics and scattering experiments, most visibly through the Green--Kubo relations that connect microscopic fluctuations to macroscopic transport coefficients.
While correlation functions are routinely computed, the particle-wise matrix from which they are assembled has received comparatively little attention as an object of study in its own right.
In this work we examine this matrix directly, drawing on random matrix theory to separate physical structure from finite-sampling noise, and on spectral information measures to track its behaviour across simulation conditions.
Using theoretical arguments together with molecular dynamics simulations of a Lennard-Jones fluid, bonded systems, water, and molten salts, we show that bonds appear as eigenvalues below the noise bulk of the particle-wise velocity correlation matrix, one per constraint, with eigenvectors that reproduce the bonding topology once the integration time exceeds the time for a pair to diffuse over its bond length, and that the von Neumann entropy of the matrix, measured against a random-matrix null, distinguishes crystal, disordered solid and liquid and locates first-order transitions at the same temperatures as the density.
